\documentclass[10pt,a4paper]{amsart}
\usepackage[margin=19mm]{geometry}
\usepackage{amsmath,amssymb,mathtools}
\usepackage[scr=rsfs,cal=euler]{mathalfa}
\usepackage{xfrac}
\usepackage[unicode,hidelinks,bookmarks=false]{hyperref}
\newcommand{\ie}{i.e.\ }
\newcommand{\cf}{cf.\ }
\newcommand{\resp}{respectively\ }
\newcommand{\Subsection}{\subsection}
\providecommand{\nobreakdash}{\nobreak\hbox{-}}
\setlength{\emergencystretch}{2em}
\multlinegap0pt
\usepackage{xcolor}
\usepackage{amssymb}
\usepackage{mathtools}
\usepackage{tikz}
\usepackage{bm}
\usepackage[shortlabels]{enumitem}
\newtheorem{proposition}{Proposition}
\newcommand{\dd}{\mathrm{d}}
\newcommand{\wtd}{\tilde d}
\newcommand{\wtdel}{\tilde\delta}
\title{Homothetic Hodge–de Rham Theory and a Geometric Regularization of Elliptic Boundary Value Problems}
\author{Fereidoun Sabetghadam}
\date{Source: arXiv:2603.27564v1 (CC BY 4.0)}
\begin{document}
\maketitle

\noindent\emph{Source and licence.} Homothetic Hodge–de Rham Theory and a Geometric Regularization of Elliptic Boundary Value Problems. Version arXiv:2603.27564v1 (CC BY 4.0), distributed by arXiv under the Creative Commons Attribution 4.0 International licence, http://creativecommons.org/licenses/by/4.0/, as recorded in the arXiv licence metadata for this exact version and shown on the abstract page https://arxiv.org/abs/2603.27564v1. Full licence notice: \texttt{licenses/arxiv-2603.27564-CC-BY-4.0.txt}.\par\smallskip

\noindent\textit{Editorial scope.} Counted unit: the article's proposition prop:twisted-adjoint (source0.tex lines 370--376). The article's complete original proof of it is delivered with it (source0.tex lines 378--397, one \texttt{proof} environment of 20 lines). No mathematics is written, repaired or re-derived: the statement and the proof are the article's own lines, in the article's own notation, and no retained line is substituted or re-flowed.  Retained uncounted as the source-local context the proof needs: lines 248--250.  Nothing was omitted from any retained span, so the package carries no omission record.  No retained line contains a citation, so the wrapper carries no bibliography and no bibliography entry.  No retained line contains a figure environment, an image-inclusion command or drawing code, so no image is delivered and the package declares no asset dependency.  Editorial formatting only, in addition: an AMS \texttt{amsart} preamble that restates the article's own declarations and macros used by the retained lines, namely the environments \texttt{proposition} on separate counters, together with the standard packages those lines need.  The retained environments print in this document as Proposition 1 in order of appearance, whereas the article prints the same items with the numbering of its own full document.  The complete native source of the article as distributed in the arXiv e-print for this exact version is bundled unchanged as \texttt{sources/197357/source0.tex}.
\begin{equation}\label{eq:weighted-pairing}
\langle \eta,\xi\rangle_{\lambda,w} \;:=\; \int_M e^{2w\lambda}\,\eta\wedge *\xi,
\end{equation}

\begin{proposition}[Twisted adjoint by conjugation]\label{prop:twisted-adjoint}
With respect to \eqref{eq:weighted-pairing}, the formal adjoint of $\wtd=e^{-w\lambda}\dd e^{w\lambda}$ is
\begin{equation}\label{eq:tilde-delta}
\wtdel \;:=\; e^{-w\lambda}\,\delta\,e^{w\lambda}.
\end{equation}
And we have $\wtdel^{\,2}=0$. 
\end{proposition}

\begin{proof}
Let $\eta\in\Omega^{p}(M)$ and $\xi\in\Omega^{p+1}(M)$ with compact support (or assume hypotheses ensuring vanishing of boundary terms, e.g.\ $M$ compact without boundary). Set $\hat\eta:=e^{w\lambda}\eta$ and $\hat\xi:=e^{w\lambda}\xi$.
Then
\[
\langle \wtd\eta,\xi\rangle_{\lambda,w}
=
\int_M e^{2w\lambda}\,(e^{-w\lambda}\dd\hat\eta)\wedge *\xi
=
\int_M (\dd\hat\eta)\wedge *\hat\xi.
\]
By definition of $\delta$ as the adjoint of $\dd$ (for the unweighted pairing), this equals
\[
\int_M \hat\eta\wedge *(\delta\hat\xi)
=
\int_M e^{2w\lambda}\,\eta\wedge *\big(e^{-w\lambda}\delta(e^{w\lambda}\xi)\big)
=
\langle \eta,\wtdel\xi\rangle_{\lambda,w}.
\]
Finally, $\wtdel^{\,2}=0$ follows from $\delta^2=0$ and the conjugation \eqref{eq:tilde-delta}.
\end{proof}

\end{document}
